\documentclass[11pt]{article}
\usepackage[a4paper,margin=22mm]{geometry}
\usepackage[T1]{fontenc}
\usepackage{lmodern}
\usepackage{amsmath}
\usepackage{microtype}
\usepackage{graphicx}
\usepackage{xcolor}
\usepackage{float}
\usepackage{fancyhdr}
\usepackage{booktabs}
\usepackage[hidelinks]{hyperref}

\definecolor{Accent}{HTML}{21918C}
\definecolor{Muted}{HTML}{5B6570}
\definecolor{Warn}{HTML}{B03A2E}

\pagestyle{fancy}\fancyhf{}
\lhead{\small\color{Muted}BayPass local-score cross-check}
\rhead{\small\color{Muted}Poikela et al., rho05 pipeline}
\cfoot{\small\thepage}
\renewcommand{\headrulewidth}{0.4pt}
\setlength{\parskip}{0.5em}\setlength{\parindent}{0pt}

\title{\vspace{-3em}\textbf{\color{Accent}BayPass's local-score outlier-region method}\\[0.2em]
\large A cross-check against the Stage-1-direct climate/mitotype scan}
\author{}
\date{\small Generated \today}

\begin{document}
\maketitle
\vspace{-2em}

\section*{\color{Accent}Motivation}
BayPass ships its own recommended method for turning per-SNP association evidence into
outlier \emph{regions}: \texttt{compute.local.scores()}
(\texttt{baypass\_public/utils/baypass\_utils.R}), implementing the local-score approach of
Fariello et al.\ (2017) and Bonhomme et al.\ (2019). This is methodologically different from the
region-definition already used throughout \texttt{module\_manuscript\_rho05} (LD-based Stage-2
(rho05) clustering, tested against a 10{,}000-draw $\Omega$-structured null covariate/contrast
distribution -- the ``floor-survivor'' test). This document asks two questions: (1) what
outlier regions does BayPass's own method identify in our four Stage-1-direct covariates
(PC1, PC2, bio\_winter, mitoC2), and (2) is that method's built-in significance threshold
actually well calibrated for this dataset, checked directly against the same
$\Omega$-structured nulls used everywhere else in this pipeline.

\textbf{Bottom line, stated up front:} the local-score method's analytic threshold is
\emph{not} well calibrated here. A meaningful fraction of pure $\Omega$-structured null draws
(no causal covariate effect at all) produce multiple ``significant'' windows using BayPass's
own default settings -- for the mitotype contrast, the typical null draw already produces
several. This document should be read as a diagnostic exercise and a cautionary
demonstration, not as a source of new candidate loci for the main manuscript.

\section*{\color{Accent}1. Method}
For each SNP, a score $S_i = \mathrm{BF(dB)}_i/10 - \xi$ (continuous covariates) or
$S_i = -\log_{10}(p)_i - \xi$ (mitoC2's C2 contrast $p$-value) is computed, where $\xi>0$ is a
fixed per-marker penalty (default $\xi=1$, used throughout). SNPs with negative BF have no
natural $p$-value, so BayPass's own code substitutes a random draw,
$-\log_{10}(U)$, $U\sim\mathrm{Unif}(0,1)$ -- this is the source of a real reproducibility
caveat discussed in Section~5. A per-chromosome \emph{Lindley process}
$L_i=\max(0, L_{i-1}+S_i)$ then accumulates evidence along the chromosome, resetting to zero
whenever cumulative evidence turns negative -- this is what lets neighbouring,
individually-weak signals combine into one region without any externally-imposed LD
clustering. An analytic significance threshold is derived per chromosome from its length and
its marker-to-marker $p$-value autocorrelation (exact for $\xi\in\{1,2,3,4\}$;
$\xi=1$ throughout here), at $\alpha=0.01$ (BayPass's default, \texttt{pval.local.score.thres}).
Markers with estimated $\hat\pi<0.2$ or $>0.8$ are excluded (\texttt{min.maf=0.2}, default).
Every excursion of $L$ above threshold is reported as one significant window
(chromosome, start, end, size, \#SNPs, and the peak position of both the Lindley score and
the original statistic within it).

Run at two resolutions:
\begin{itemize}
  \item \textbf{Full-SNP} (\texttt{local\_score\_regions.R}): every genome-wide SNP
    individually ($\sim$1.08M retained after MAF filtering per covariate), using the real
    full-SNP BayPass scan already computed for the manuscript.
  \item \textbf{Stage-1-cluster} (\texttt{local\_score\_regions\_S1units.R}): the 18{,}361
    Stage-1-direct clusters ($n_{\mathrm{SNPs}}\geq5$) used by Module~C, one representative
    marker per cluster. \texttt{min.nsnp} (minimum markers/chromosome to compute a threshold)
    had to be lowered from BayPass's default (10{,}000) to 100 -- the largest chromosome has
    only 1{,}078 Stage-1 units, so the default would silently exclude every chromosome.
\end{itemize}

\section*{\color{Accent}2. Observed data, full-SNP resolution}
\begin{table}[H]\centering\small
\begin{tabular}{lrrl}
\toprule
Covariate & Raw threshold crossings & Local-score windows & Total window span \\
\midrule
PC1        & 2{,}779 & 6  & Chr2, Chr8, Chr15 ($\times2$), Chr24 ($\times2$) \\
PC2        & 1{,}484 & 7  & Chr4, Chr5, Chr13, Chr24, Chr26 ($\times2$), Chr27 \\
bio\_winter & 2{,}636 & 8  & Chr2, Chr3, Chr8, Chr11, Chr15 ($\times2$), Chr24 ($\times2$) \\
mitoC2      & 648     & 95 & scattered across 26 of 26 chromosomes \\
\bottomrule
\end{tabular}
\caption{Local-score windows, full-SNP resolution, BayPass default parameters
($\xi=1$, $\alpha=0.01$, MAF$\geq0.2$). Figure~\ref{fig:obs-fullsnp}.}
\end{table}

\textbf{Striking cross-covariate overlap.} PC1 and bio\_winter's windows share
near-\emph{identical} start coordinates at five of their six/eight loci (Chr2: 10{,}008{,}320
bp both; Chr8: 6{,}712{,}279 bp both; Chr15: both peaks within $\sim$1.2~kb; Chr24: both peaks
within $\sim$0 bp). PC2 additionally has its own Chr24 window (1{,}764{,}859--1{,}843{,}054)
\emph{nested entirely inside} PC1/bio\_winter's broader Chr24 window
(1{,}309{,}506--1{,}894{,}413/1{,}928{,}191). Three unrelated covariates converging on the same
loci is not consistent with covariate-specific climate signal -- it is exactly the
genome-architecture/ancestry-confound pattern already central to
\texttt{ancestry\_climate\_mitotype.tex} (Section~6 there), now recovered independently by a
completely different region-detection method.

\textbf{Relationship to the floor-survivor candidates.} Of the 12 floor-survivor loci
identified elsewhere in this pipeline (PC1: 1, PC2: 2, bio\_winter: 10), only bio\_winter's
three loci on Chr2 (both units) and Chr24 fall inside a full-SNP local-score window; PC1's
Chr24 floor-survivor sits in the $\sim$170~kb gap between two adjacent local-score windows,
and neither PC2 floor-survivor is recovered at all. The two methods substantially disagree at
the level of individual candidate loci, even though (per the paragraph above) they agree on
which broad regions show elevated signal shared across covariates.

\begin{figure}[H]\centering
  \includegraphics[width=0.9\linewidth]{../Figures/local_score_snp_manhattan_combined.png}
  \caption{\label{fig:obs-fullsnp}Full-SNP local-score windows, styled identically to the
  manuscript's Figure~5 (\texttt{ancestry\_climate\_mitotype.tex}): every genome-wide SNP at
  its own BayPass value; grey = not in any raw-crossing Stage-2 (rho05) region; coloured =
  Stage-2 region containing $\geq1$ raw threshold crossing (BF(dB)$\geq15$, or
  $-\log_{10}(p)\geq3$ for mitoC2); arrowed and labelled = one arrow per significant
  local-score window (not per touched Stage-2 region -- see Section~4), positioned at the
  window midpoint and labelled/coloured by the Stage-2 group of the window's own peak SNP.}
\end{figure}

\section*{\color{Accent}3. Observed data, Stage-1-cluster resolution}
\begin{table}[H]\centering\small
\begin{tabular}{lrrp{7cm}}
\toprule
Covariate & Local-score windows & vs. full-SNP & Windows (chr) \\
\midrule
PC1        & 1 & 6 & Chr24 \\
PC2        & 0 & 7 & -- \\
bio\_winter & 2 & 8 & Chr2, Chr24 \\
mitoC2      & 8 & 95 & scattered \\
\bottomrule
\end{tabular}
\caption{Local-score windows, Stage-1-cluster resolution (18{,}361 units, min.nsnp=100).
Figure~\ref{fig:obs-s1}.}
\end{table}

Far fewer windows survive at this coarser resolution -- expected, since Stage-1 units are
deliberately LD-pruned (rho=0.5) to be much less correlated than neighbouring full SNPs, so
the Lindley process has far less to accumulate. PC1's Chr24 window and bio\_winter's Chr2 and
Chr24 windows correspond to the SAME loci recovered at full-SNP resolution and to real
floor-survivor candidates -- the one part of this cross-check where local-score, Stage-1, and
the structured-null floor test all agree. PC2 loses its signal entirely at this resolution.

\begin{figure}[H]\centering
  \includegraphics[width=0.78\linewidth]{../Figures/local_score_manhattan_S1units.png}
  \caption{\label{fig:obs-s1}Local-score windows, Stage-1-cluster resolution. Points are
  coloured by local-score window identity (not Stage-2 group, unlike
  Figure~\ref{fig:obs-fullsnp}); every genome-wide member SNP of a tested Stage-1 cluster
  inherits its cluster's own single value (not an independent per-SNP estimate -- no such
  Stage-1-resolution SNP-level scan exists), which is why points inside a coloured window form
  flat vertical bands rather than the natural scatter in Figure~\ref{fig:obs-fullsnp}.}
\end{figure}

\clearpage
\section*{\color{Accent}4. Does the analytic threshold survive the $\Omega$-structured null?}
The real question this document exists to answer: if a covariate/contrast carries
\emph{no} causal signal -- only the populations' $\Omega$ covariance structure, exactly like
the null covariates already built for Module~C's genome-wide calibration
(\texttt{moduleC\_stage1\_null\_regen.R}) and the mitotype floor test
(\texttt{moduleB\_stage1\_mitoC2\_null.R}) -- how often does BayPass's own analytic threshold
still call something ``significant''? Run at Stage-1-cluster resolution (the only resolution
at which null draws exist at all -- generating a full-SNP null would need $\Omega$-structured
covariates re-run through BayPass on $\sim$1.1M markers each, an order of magnitude beyond
anything already computed in this pipeline).

\begin{table}[H]\centering\small
\begin{tabular}{lrrrrr}
\toprule
Null pool (1{,}000 draws each) & Mean windows & Median & Max & \% draws $\geq$ real count \\
\midrule
Continuous (PC1/PC2/bio\_winter's shared null) & 0.25 & 0 & 3 & 21.6\% ($\geq1$) \\
mitoC2's dedicated contrast null                & 3.88 & 4 & 16 & 6.6\% ($\geq8$) \\
\bottomrule
\end{tabular}
\caption{Distribution of local-score window counts under 1{,}000 structured-null draws per
pool (\texttt{local\_score\_null\_check\_S1units.R}), same procedure as the observed data.
``\% draws $\geq$ real count'' compares against the real covariates' OWN Stage-1-cluster
window count from Section~3 (PC1: 1, bio\_winter: 2, mitoC2: 8 -- PC2's real count is 0, so no
comparable tail probability applies).}
\end{table}

\textbf{Interpretation.} Roughly one in five draws from a covariate with \emph{zero} true
signal already produces at least one ``significant'' local-score window under BayPass's
default threshold; for the mitotype contrast, the typical null draw produces about four. This
is a strong, direct demonstration that the analytic threshold under-corrects for this
dataset's population structure -- consistent with, and now quantifying, the same conclusion
reached independently in \texttt{ancestry\_climate\_mitotype.tex} Section~6 (mitoC2's
genome-wide BF/C2 landscape is indistinguishable from its own $\Omega$-structured null on
every measure tested there).

\section*{\color{Accent}5. False-positive demonstration}
Figure~\ref{fig:null-demo} renders one high-scoring null draw per covariate --
\emph{styled identically to Figure~\ref{fig:obs-fullsnp}} (Stage-2 (rho05) region colouring,
one arrow per significant window) -- to make concrete how convincing a purely-null result can
look. PC1/PC2/bio\_winter each use a different draw from the shared continuous-null pool
(\#624, \#536, \#107: 3, 2, 2 windows) and mitoC2 uses its own pool's worst draw (\#480: 16
windows). None of these draws carry any real covariate signal.

\begin{figure}[H]\centering
  \includegraphics[width=0.78\linewidth]{../Figures/local_score_null_examples.png}
  \caption{\label{fig:null-demo}Four \textbf{NULL} draws (titled in red, explicitly marked
  NOT REAL DATA), one per covariate, styled identically to Figure~\ref{fig:obs-fullsnp}.
  Every arrowed, labelled, distinctly-coloured region here is a false positive by
  construction.}
\end{figure}

\section*{\color{Accent}6. Caveats and limitations}\label{sec:limitations}
\begin{itemize}
  \item \textbf{The BF-based score is not deterministic.} BayPass's own code draws a random
    $p$-value for every negative BF value (Section~1), so re-running
    \texttt{compute.local.scores()} on the identical BF vector can change the exact window
    count and boundaries from call to call. All numbers in this document use a fixed seed
    (\texttt{set.seed(1)} in the region-calling scripts; \texttt{set.seed(10000+draw)} per
    panel in the demonstration figure) purely for reproducibility -- the instability itself is
    a genuine, separate calibration concern for continuous covariates, on top of Section~4's
    finding. mitoC2 (a $p$-value statistic) is unaffected and fully deterministic.
  \item \textbf{Arrows are per window, not per touched Stage-2 (rho05) region.} An early
    version of the full-SNP figure arrowed every Stage-2 cluster overlapped by a significant
    window and produced 138--8{,}055 arrows per covariate, because local-score windows
    (hundreds of kb to several Mb) routinely span dozens of the much finer-grained
    (rho=0.5, 0.5~cM cap) Stage-2 clusters. That inflated the apparent candidate count and
    does not reflect what the method actually claims (a WINDOW is the significant unit).
    Every figure in this document instead places one arrow per window, at its midpoint,
    labelled/coloured by the Stage-2 group of the window's own peak SNP -- which is itself
    often not a raw threshold crossing (the method's entire point is windows where evidence
    accumulates without any single SNP crossing a naive threshold).
  \item \textbf{No full-SNP-resolution null exists.} Section~4/5's calibration check and
    demonstration are necessarily at Stage-1-cluster resolution. The qualitative conclusion
    (the analytic threshold under-corrects for this dataset) is unlikely to reverse at full-SNP
    resolution -- if anything, denser, more autocorrelated markers give the Lindley process
    more opportunity to accumulate spurious runs -- but this has not been directly tested and
    would require a large new BayPass campaign to check properly.
  \item \textbf{This document does not replace the floor-survivor test.} The
    $\Omega$-structured-null-calibrated floor-survivor test used throughout
    \texttt{module\_manuscript\_rho05} remains the primary, trusted candidate-locus method for
    this project. Local-score is presented here purely as an independent cross-check and a
    calibration diagnostic, not as a source of new candidates for the manuscript.
\end{itemize}

\section*{\color{Accent}Summary}
BayPass's own local-score outlier-region method, run with its own default settings, recovers
a striking cross-covariate architecture signal (PC1, PC2, and bio\_winter converging on
near-identical genomic windows despite carrying no obvious shared cause) that reinforces the
manuscript's ancestry-confound narrative. But a direct check against the same
$\Omega$-structured nulls used everywhere else in this pipeline shows the method's analytic
significance threshold is not well calibrated for this dataset: roughly one in five
structure-only null draws for the continuous covariates, and the clear majority for mitoC2's
contrast, already produce one or more ``significant'' windows by chance. Combined with a
separate, unrelated instability in how the method scores negative Bayes factors, this method
should be treated as a supplementary, exploratory cross-check -- informative for the
architecture question it incidentally answers, but not a substitute for, or independent
confirmation of, the structured-null floor-survivor candidates already reported in the main
manuscript.

\bigskip
{\small\color{Muted}
\textit{References:} Fariello M.-I. et al.\ (2017) \textit{Mol.\ Ecol.} 26:3700--3714 --
local score method. Bonhomme M.\ et al.\ (2019) \textit{Heredity} 123:1--14 -- extension to
$\xi=3,4$. Olazcuaga L.\ et al.\ (2020) -- BayPass C2 contrast statistic.}

\end{document}
